\subsection{Ausgangsituation}\label{subsec:Ausgangssituation}

\par Die \gls{stz} wurde 2023 gegründet und dient der Verwaltung und Bewirtschaftung von Immobilien, die von der \gls{stz} erworben werden. Gleichwohl die \gls{kg}-Gründung erst 2023 erfolgte, befanden sich die in der \gls{kg} gebündelten Gesellschaften bereits teilweise jahrzehntelang im Besitz und unter der Verwaltung des Kompelemtärs und der Kommanditisten der \gls{stz}. Das Ziel der jungen \gls{kg} ist die nachhaltige Sicherung des bestehenden Immobilienportfolios samt Übertragung auf die jüngeren Gesellschafter, das bestehende Immobilienportfolio zu erweitern und ein attraktives Immobilienportfolio aufzubauen und zu verwalten, um langfristige Wertsteigerungen zu erzielen. Die \gls{stz} verfolgt eine nachhaltige Anlagestrategie und legt dabei besonderen Wert auf die Auswahl von Immobilien in
attraktiven Lagen mit einem hohen Entwicklungspotenzial.

\par Die \gls{stz} ist eine familiengeführte \gls{kg} bestehen aus einem Komplementär, der die Geschäftsführung übernimmt, und mehreren Kommanditisten, die als Kapitalgeber fungieren.

\begin{itemize}
    \item Komplementär: Gotthard Seitz (06.07.1957)
    \item Kommanditistin: Angela Seitz (01.08.1963)
    \item Kommanditist: Sebastian Seitz (12.11.1987)
    \item Kommanditist: Fabian Seitz (25.04.1991)
\end{itemize}

\par Im besitz und Verwaltung der \gls{stz} befinden sich derzeit mehrere Immobilien, darunter Wohn- und Gewerbeimmobilien in verschiedenen Städten. Die \gls{stz} plant, ihr Portfolio durch den Erwerb weiterer Immobilien zu erweitern, insbesondere in München, um von der positiven Entwicklung des Immobilienmarktes in der Region zu profitieren. Die \gls{stz} verfolgt dabei eine langfristige Anlagestrategie und legt besonderen Wert auf die Auswahl von Immobilien in attraktiven Lagen mit einem hohen Entwicklungspotenzial.